\documentclass[a4paper,12pt,twoside]{article}

\usepackage[german,english]{babel}

%%% FOR DISPLAYING CODE
\usepackage[procnames]{listings}
\usepackage{xcolor}

%%% MARGINS
\usepackage{geometry}
% INFO: https://www.overleaf.com/learn/latex/Page_size_and_margins
\geometry{
	a4paper,
%	total={170mm,257mm},
	left=25mm,
	top=25mm,
	textwidth=16cm,
	textheight=23cm,
%	right=25mm,
%	bottom=40mm,
}

%%% SPACING
%\usepackage{setspace} 
%\onehalfspacing

\usepackage{siunitx} % for writing numbers with units % e.g. \SI{6.3e8}{km}, \ang{62.59}, \num{100}'s \si{km}, \SI{53}{\degree}
\usepackage{amsmath,amssymb}
\usepackage{graphicx}
\usepackage{datetime} 
\usepackage{color}
\usepackage{float} %without it, tables position more or less random, with this an writing \begin{table}[H] the table appears right where it is defined!
\usepackage{mathptmx}
\usepackage{enumitem}
\usepackage[T1]{fontenc}

\usepackage[utf8]{inputenc}

%%% HEADER and FOOTER
%\usepackage[automark,headsepline]{scrpage2}
%\ihead[]{\headmark}
%\chead[]{}
%\ohead[]{\pagemark}
%\ifoot[]{}
%\cfoot[\pagemark]{}
%\ofoot[]{}
%\pagestyle{scrheadings}

\usepackage{lipsum}

%% LINE SPACING
\usepackage{setspace}
\setstretch{1.5}

%%% HYPERLINKS
\usepackage{hyperref} % muss sich vor cite stuff befinden

%%% CITE
\usepackage{apacite}
\bibliographystyle{apacite}

%!TEX root = ../maturaarbeit.tex

%%%%%%%%%%%%%%%%%%%%%%%%%%%%%%%%%%
% DEFINE COLORS
\definecolor{red}{rgb}{0.6,0,0} 
\definecolor{blue}{rgb}{0,0,0.6}
\definecolor{green}{rgb}{0,0.8,0}
\definecolor{cyan}{rgb}{0.0,0.6,0.6}

\definecolor{mypink}{rgb}{0.753,0.000,0.890}
\definecolor{myblue}{rgb}{0.078,0.000,1.000}
\definecolor{mybluedark}{rgb}{0.004,0.024,0.525} \definecolor{mygreen}{rgb}{0.000,0.514,0.000}
\definecolor{myreddark}{rgb}{0.698,0.000,0.008}
\definecolor{mycyan}{rgb}{0.000,0.506,0.612}
\definecolor{mybrown}{rgb}{0.494,0.365,0.090}

%%% DISPLAY CODE
\usepackage{listings}
\newcommand\pythonstyle{\lstset{
    language=Python,
	tabsize=4,
	basicstyle=\normalsize\sffamily,
	numberstyle=\color{gray},
	stringstyle=\color{myreddark},
    commentstyle=\color{mygreen},
    % KEYWORDS
    % main keywords
	keywordstyle=\normalsize\color{myblue},%\bfseries,
    % add keywords (main blue)
    emph={False,None,True,self,TODO},
    emphstyle={\color{myblue}},
    % pink emph
    emph={[2]assert,break,continue,del,elif ,else,except,finally,for,from,global,if,import,in,pass,raise,return,try,while,with,yield},
    emphstyle={[2]\color{mypink}},%\bfseries,
    %dark blue emph
    emph={[3]execfile,reduce,xrange},
    emphstyle={[3]\color{mybluedark}},
    % brown emph
    emph={[4]exec,print,isinstance,zip,enumerate,reversed,len,repr},
    emphstyle={[4]\color{mybrown}},
    % cyan emph
    emph={[5]object,type,list,set,dict,tuple,str,super},
    emphstyle={[5]\color{mycyan}},
    % errors (also cyan emph)
    emph={[6]Exception,NameError,IndexError,SyntaxError,TypeError,ValueError,OverflowError,ZeroDivisionError},
    emphstyle={[6]\color{mycyan}},
    % errors (also cyan emph)
    emph={[7]copy,deepcopy,append,real,imag},
    emphstyle={[7]\color{black}},
    % 
    showstringspaces=false,
	breaklines=true,
	numbers=left,
    frame=tb,
	xleftmargin=15pt
}}

\newcommand\javascriptstyle{\lstset{
	basicstyle=\small\sffamily,
    keywordstyle=\color{blue}\bfseries,
    commentstyle=\color{purple}\ttfamily,
    stringstyle=\color{red}\ttfamily,
	numberstyle=\color{gray},
    keywords={typeof, new, true, false, catch, function, return, null, catch, switch, var, if, in, while, do, else, case, break},
    ndkeywords={class, export, boolean, throw, implements, import, this},
    ndkeywordstyle=\color{darkgray}\bfseries,
    identifierstyle=\color{black},
    comment=[l]{//},
    morecomment=[s]{/*}{*/},
    morestring=[b]',
    morestring=[b]",
    showstringspaces=false,
	breaklines=true,
	numbers=left,
    frame=tb,
	xleftmargin=15pt,
    sensitive=false,
}}

\newcommand\csharpstyle{\lstset{
	language=csh,
	tabsize=4,
	basicstyle=\small\sffamily,
	numberstyle=\color{gray},
	stringstyle=\color{myreddark},
    commentstyle=\color{mygreen},
	morecomment=[l]{//}, %use comment-line-style!
	morecomment=[s]{/*}{*/}, %for multiline comments
    % KEYWORDS
	keywordstyle=\normalsize\color{myblue},%\bfseries,
	morekeywords={ abstract, event, new, struct,
		as, explicit, null, switch,
		base, extern, object, this,
		bool, false, operator, throw,
		break, finally, out, true,
		byte, fixed, override, try,
		case, float, params, typeof,
		catch, for, private, uint,
		char, foreach, protected, ulong,
		checked, goto, public, unchecked,
		class, if, readonly, unsafe,
		const, implicit, ref, ushort,
		continue, in, return, using,
		decimal, int, sbyte, virtual,
		default, interface, sealed, volatile,
		delegate, internal, short, void,
		do, is, sizeof, while,
		double, lock, stackalloc,
		else, long, static,
		enum, namespace, string},
	% 
    showstringspaces=false,
	breaklines=true,
	numbers=left,
    frame=tb,
	xleftmargin=15pt	
}}

% Python environment
\lstnewenvironment{python}[1][]
{
	\pythonstyle
	\lstset{#1}
}
{}
\lstnewenvironment{javascript}[1][]
{
	\javascriptstyle
	\lstset{#1}
}
{}
\lstnewenvironment{csharp}[1][]
{
	\csharpstyle
	\lstset{#1}
}
{}

% CODE FOR EXTERNAL FILES
\newcommand\pythonexternal[2][]{{
		\pythonstyle
		\lstinputlisting[#1]{#2}}}

% CODE FOR INLINE
\newcommand\pythoninline[1]{{\pythonstyle\lstinline!#1!}}
\newcommand\csharpinline[1]{{\csharpstyle\lstinline!#1!}}
\newcommand\javascriptinline[1]{{\javascriptstyle\lstinline!#1!}}

%% DEFINE CUSTOM COMMANDS AND SHORTCUTS

% some commands
\def\ba#1\ea{\begin{align}#1\end{align}}
\def\bas#1\eas{\begin{align*}#1\end{align*}}
\def\bmat#1\emat{\begin{pmatrix}#1\end{pmatrix}}
\newcommand{\ve}[1]{\vec{#1}}
\newcommand{\veTwo}[2]{\begin{pmatrix}#1\\#2\end{pmatrix}}
\newcommand{\veThree}[3]{\begin{pmatrix}#1\\#2\\#3\end{pmatrix}}
\newcommand{\veFour}[4]{\begin{pmatrix}#1\\#2\\#3\\#4\end{pmatrix}}
\newcommand{\ora}[1]{\overrightarrow{#1}}
\newcommand{\ola}[1]{\overleftarrow{#1}}
\newcommand{\s}[1]{\sqrt{#1}}

\def \bit{\begin{itemize}\setlength\itemsep{0em}} %\vspace{-5mm}
	\def \eit{\end{itemize}}

\def \ben{\begin{enumerate}\setlength\itemsep{0em}} %\vspace{-5mm}
	\def \een{\end{enumerate}}

\def \N{\mathbb{N}}
\def \Z{\mathbb{Z}}
\def \Q{\mathbb{Q}}
\def \R{\mathbb{R}}
\def \C{\mathbb{C}}

\def \ra{\rightarrow}
\def \longra{\longrightarrow}
\def \Ra{\Rightarrow}
\def \Longra{\Longrightarrow}
\def \la{\leftarrow}
\def \longla{\longleftarrow}
\def \La{\Leftarrow}
\def \Longla{\Longleftarrow}

\def \lra{\leftrightarrow}
\def \longlra{\longleftrightarrow}
\def \Lra{\Leftrightarrow}
\def \Longlra{\Longleftrightarrow}

\def \l{\left}
\def \r{\right}

\def \a{\alpha}
\def \b{\beta}
\def \g{\gamma}

\def \c{\cdot}

\def \el{\in}
\def \notel{\notin}

\newcommand{\f}{\frac}

\def \q{\quad}
\newcommand{\p}{\phantom}

\def\bs#1\es{\begin{split}#1\end{split}}


\def \L{\mathbb{L}}




\begin{document}

\selectlanguage{german} %% HIER SPRACHE EINSTELLEN! english, german

%% TITEL PAGE
%\title{Titel des Dokumentes}
%\author{Max Muster}
%\date{01.01.0001}
%\maketitle

%!TEX root = ../maturaarbeit.tex

\begin{titlepage}
	\clearpage\thispagestyle{empty}	
	\setstretch{1}
	
	\begin{minipage}[t]{\textwidth}
		\begin{minipage}[t]{0.5\textwidth}
			Hanna Kradolfer\\
			Götighoferstrasse 11\\
			8586 Riedt b. Erlen\\
			077 479 79 01\\
			hanna.kradolfer@gmail.com
		\end{minipage}
		\begin{minipage}[t]{0.5\textwidth}
			\begin{flushright}
				Kantonsschule Romanshorn\\
				Klasse 3mdz\\
				SLA
			\end{flushright}
		\end{minipage}
	\end{minipage}
	
	\vspace{4cm}
	
	{
		\centering
		\Huge\bfseries Einfluss dunkler Materie aud die Rotationskurve von Galaxien\par
		\vspace{1cm}
		\includegraphics[width=0.15\textwidth]{example-image-1x1}\par
	}
	
	\vspace{9cm}	
	\noindent
	Fach: Physik \noindent\\
	Betreuungsperson: Dr. Andreas Schärer\\
	Abgabetermin: 11.9.2023
	
\end{titlepage}


% \null\thispagestyle{empty}\clearpage

\pagenumbering{roman}

%!TEX root = ../maturaarbeit.tex

\section*{Abstract}

Die Messungen und Berechnungen der Rotationskurven von Galaxien verdeutlichen, dass Dunkle Materie einen massgeblichen Einfluss auf diese Kurven ausüben muss. Die beobachteten Konstanten oder ansteigenden Geschwindigkeiten in den äusseren Galaxienregionen widersprechen den Erwartungen, die alleine auf sichtbarer Materie basieren. Diese Erkenntnis belegt die entscheidende Rolle der Dunklen Materie bei der Erklärung der Rotationsdynamik von Galaxien.

%\null\thispagestyle{empty}\clearpage

%\null\thispagestyle{empty}\clearpage

\newpage

%% TABLE OF CONTENTS
%\setcounter{page}{1}
\tableofcontents

\parindent=0pt
\parskip=6pt

\newpage

\pagenumbering{arabic}

%%%%%%%% DEIN CODE AB HIER

%!TEX root = ../maturaarbeit.tex

\section{Einleitung}

E sh ihleitig


\newpage
%!TEX root = ../maturaarbeit.tex


\section{Die vier fundamentalen Wechselwirkungen}

Die vier Wechselwirkungen setzen sich aus der Gravitation, der elektromagnetischen Wechselwirkung, der schwachen Wechselwirkung und der starken Wechselwirkung zusammen. Jede Wechselwirkung hat eine eigene Ladung. Es können alle Phänomene und Prozesse, welche auf der Erde oder im Weltall beobachtet werden, mit den vier fundamentalen Wechselwirkungen beschrieben werden.
Das Standardmodell der Teilchenphysik beschreibt drei der vier Wechselwirkungen, die Gravitation spielt keine zentrale Rolle, da die Teilchen eine sehr kleine Masse haben. Die Gravitation wird durch die allgemeine Relativitätstheorie beschrieben, die nicht Teil des Standardmodells ist. Sie wirkt zwischen allen Teilchen, welche eine Masse besitzen. 
Die elektromagnetische Wechselwirkung wirkt zwischen elektrisch geladenen Teilchen wie Elektronen und Protonen. Die schwache Wechselwirkung wirkt bei der Kernfusion. Die starke Wechselwirkung hält die Protonen zusammen.
Die Gravitation und die elektromagnetische Wechselwirkung erfahren wir direkt in unserem Alltag, die schwache und starke Wechselwirkung jedoch nicht, da ihre Reichweite zu gering ist \cite{Bührke2022}.

\section{Dunkle Materie}

Es gibt zwei Arten von Dunkler Materie, die baryonische und die nicht baryonische Dunkle Materie.\\ Unter baryonischer Materie versteht man normale Materie, welche aus Elektronen, Neutronen und Protonen besteht. Damit sind zum Beispiel Objekte wie massenarme und daher leuchtschwache Sterne gemeint. Diese haben die Fachbezeichnung \glqq MACHO\grqq{} was für \glqq MAssive Compact Halo Objects\grqq{} steht. Diese können nicht direkt beobachtet werden. Braune Zwerge, Weisse Zwerge und Neutronensterne sowie Schwarze Löcher gehören auch zu der Kategorie baryonische Dunkle Materie \cite{Bührke2022}.
Jedoch reicht die Masse dieser Objekte nicht aus, um alle Beobachtungen im All zu erklären\cite{Moeller2010}.\\ Nicht baryonische Materie, wie der Name verrät, besteht nicht aus Elektronen, Neutronen und Protonen. Es gibt verschiedene Vermutungen, was diese sein könnten \cite{Anderl2023}.\\
In Frage kommen WIMPs (weakly interacting massive particles), was auf Deutsch schwach wechselwirkende, massenreiche Teilchen bedeutet. Diese Teilchen besitzen keine Ladung und somit auch kein elektrisches oder magnetisches Feld. Dadurch beschränkt sich ihre Wechselwirkung auf die Gravitation und die schwache Wechselwirkung. Das heisst, dass Dunkle Materie ungestört Planeten durchqueren könnte, wodurch sie auch so schwer zu detektieren sind. 
Es wird mit verschiedenen Methoden nach WIMPs gesucht. Die Suchmethoden basieren darauf, dass ein WIMP mit einem Atomkern zusammenstösst und das kann unterschiedliche Folgen haben, wie die Ionisation, die Szintillation oder die Phononen-Anregung. Bei der Ionisation schlägt das WIMP ein Elektron aus der Atomhülle heraus und ein messbarer elektrischer Strom wird erzeugt.
Bei der Szintillation löst das WIMP einen kurzen Lichtblitz aus, wenn es auf einen Atomkern prallt. Bei der Phononen-Anregung kann die Kollision mit einem WIMP in einem Kristall Schwingungen des Kristallgitters auslösen, und dabei erhöht sich die Temperatur geringfügig. Jedoch wurde bis heute die Existenz von WIMPs noch nicht nachgewiesen. \cite{Bührke2022} 

\newpage
%!TEX root = ../maturaarbeit.tex

\section{Kapitel 3}

\lipsum[4]

\subsection{Baryonische Materie}

Füge ein Bild ein:
\begin{figure}[H]
	\centering
	\includegraphics[width=0.4\textwidth]{figures/grumpy_cat}
	\caption[Eintrag in Abbildungsverzeichnis von Grumpy Cat]{Grumpy cat. The Grumpy Cat is a grumpy cat.}
	\label{fig grumpy cat.}
\end{figure}
Natürlich kann man auch Referenzen auf Bilder einfügen.
Bild \ref{fig grumpy cat.} zeigt die Grumpy Cat.
Die Breite/Höhe eines Bildes kann man angeben z.B. mit width=3cm oder height=15mm.
Soll das Bild die halbe Breite des Textes haben, kann man angeben
width=0.5\textbackslash textwidth

\newpage

\subsection{Dunkle Materie}

Liste
\begin{itemize}
	\item Erstes Element
	\item Zweites Element
	\item Drittes Element
	\item Viertes Element
\end{itemize}

Sehr kompakte Liste:
\begin{itemize}
	\vspace{-\topsep}
	\setlength{\itemsep}{0pt}\setlength{\parskip}{0pt}
	\item Erstes Element
	\item Zweites Element
	\item Drittes Element
	\item Viertes Element
\end{itemize}

\subsection{Aufzählung}

Aufzählung mit Zahlen nummeriert:
\begin{enumerate}
	\item Erstes Element
	\item Zweites Element
	\item Drittes Element
	\item Viertes Element
\end{enumerate}

Aufzählung mit Grossbuchstaben nummeriert:
\begin{enumerate}[label=\Alph*]
	\item Erstes Element
	\item Zweites Element
	\item Drittes Element
	\item Viertes Element
\end{enumerate}

Aufzählung mit krossbuchstaben nummeriert:
\begin{enumerate}[label=\alph*)]
	\item Erstes Element
	\item Zweites Element
	\item Drittes Element
	\item Viertes Element
\end{enumerate}

Aufzählung mit römischen Buchstaben nummeriert:
\begin{enumerate}[label=(\roman*)]
	\item Erstes Element
	\item Zweites Element
	\item Drittes Element
	\item Viertes Element
\end{enumerate}

\subsection*{Kapitel ohne Nummerierung}

Wer meine Nummer findet, soll sich bitte bei mir melden!

\subsection{Hilfreiche Links}

\subsubsection{Tabelle}

Verwende den folgenden Tabellen-Generator, um einfache Tabellen zu erzeugen: \url{https://www.tablesgenerator.com}

\subsubsection{Symbole}

Falls du nicht weisst, wie man ein gewisses Symbol in LaTeX darstellst, so hilft dir folgende Seite, um dieses zu identifizieren: \url{http://detexify.kirelabs.org/classify.html}

% \lipsum[1]

\subsection{Zweites Unterkapitel des zweiten Kapitels}

\lipsum[2-3]

\newpage
%!TEX root = ../maturaarbeit.tex

\section{Schlusswort}

Die Untersuchung der Rotationskurven von Galaxien hat gezeigt, dass die beobachteten Geschwindigkeiten der Sterne und der Gase in den äusseren Regionen von Galaxien oft nicht den Erwartungen entsprechen, die auf Grundlage der sichtbaren Materie allein getroffen wurden. Diese Diskrepanz zwischen den beobachteten Rotationskurven und den theoretisch vorhergesagten Kurven kann nur durch die Anwesenheit von Dunkler Materie erklärt werden.
Dunkle Materie, obwohl unsichtbar und noch nicht bewiesen, übt eine starke Gravitationskraft aus, die die sichtbare Materie in einer Galaxie beeinflusst. Die Dunkle Materie bildet eine unsichtbare Halo-Struktur um das Zentrum der Galaxie herum und trägt wesentlich zur Rotationsdynamik bei. Die Anwesenheit von Dunkler Materie erklärt, warum die Rotationskurven in den äusseren Regionen von Galaxien flach bleiben, anstatt wie erwartet, abzufallen.

Obwohl Dunkle Materie noch nie nachgewiesen wurde, deutet sehr viel darauf hin, dass es sie gibt. Mich fasziniert, wie durch Beobachten von Sternen, die über Tausende von Lichtjahren entfernt sind, so viel herausgefunden werden kann.
Der Gedanke, dass Dunkle Materie ein so wichtiger Teil sein muss, finde ich erstaunlich.
 








\newpage
\appendix
% %!TEX root = ../maturaarbeit.tex

\section{Erster Anhang}

Ergänzende Informationen gehören in den Anhang. Wären diese Informationen im Haupttext der Arbeit, würden sie stören, zum Beispiel weil es zu viel und zu detailliert ist. In den Anhang (unter Umständen) gehören:
\begin{itemize}
	\item Code einer Programmierarbeit (Zeige im Haupttext nur Codeausschnitte, die der Argumentation helfen)
	\item detaillierte Berechnungen
	\item Ergebnisse einer Umfrage      
	\item Interviews
	\item \ldots
\end{itemize}


% \newpage
% %!TEX root = ../maturaarbeit.tex

\section{Zweiter Anhang}

Anhang der 2.


%%%%%%%% DEIN CODE BIS HIER

\newpage

\clearpage

\bibliography{tex_content/biblio}

\newpage
\addcontentsline{toc}{section}{\listfigurename} % Fuege Abb.vz zu TOC hinzu
\listoffigures % Abbildungsverzeichnis
\newpage
%!TEX root = ../maturaarbeit.tex



Hiermit versichere ich, dass ich die vorstehende Arbeit selbstständig angefertigt habe.
Alle Stellen, die wörtlich oder sinngemäss übernommen wurden, habe ich als solche
gekennzeichnet. 

Hanna Kradolfer


\end{document}

